\clearpage
\section{Hearing Harmonics}\label{sec:harmonics}

A spectrum connects a sound's timbre to its partials. Try this with an
oscillator that offers sine, triangle, saw, and square outputs at one pitch.

\paragraph{Starting point}
Load \textbf{Harmonics}. Set all input gains to $0\,\mathrm{dB}$, Length
4096, Hop $30\,\mathrm{ms}$, and Y Scale to Log 60dB. Keep Average and
Smooth off, Slope at zero, and the view at $0$--$5\,\mathrm{kHz}$.
Tune the oscillator near $220\,\mathrm{Hz}$; exact tuning is not required.

\begin{enumerate}
  \item Patch sine to red and saw to green. Look for a shared lowest peak,
  the fundamental. The saw should have peaks near integer multiples of it:
  $440$, $660$, $880\,\mathrm{Hz}$, and so on.
  \item Patch triangle to blue and a square with 50\% pulse width to yellow.
  Compare peaks near $3f_0$ and $5f_0$. An ideal triangle's upper harmonics
  fall faster than a square's; both lack even harmonics.
  \item Change the pulse width. Even harmonics can now appear, and some
  harmonics weaken or disappear. Listen while watching the pattern change.
  \item Increase Smooth, then return it to None. The broad tonal shape is
  easier to compare when smoothed, but individual partials merge.
\end{enumerate}

\begin{figure}[!htp]
\centering
% Ideal relative partial amplitudes, independent of module UI and FFT results.
\begin{tikzpicture}[x=1cm,y=1cm,font=\small]
\foreach \origin/\title/\mode in {0/Saw/1,4.5/Square/2,9/Triangle/3} {
  \begin{scope}[shift={(\origin,0)}]
    \node[anchor=west,font=\small\bfseries,text=fourierAccent] at (0,2.25) {\title};
    \draw[->,panelMuted] (0,0)--(3.85,0);
    \draw[->,panelMuted] (0,0)--(0,1.85);
    \foreach \n in {1,...,7} {
      \pgfmathsetmacro{\amp}{\mode==1 ? 1/\n : (mod(\n,2)==0 ? 0 : (\mode==2 ? 1/\n : 1/(\n*\n)))}
      \draw[fourierAccent,line width=1.2pt] (\n*.45,0)--(\n*.45,\amp*1.65);
      \fill[fourierAccent] (\n*.45,\amp*1.65) circle[radius=.035];
    }
    \node[below] at (.45,0) {$f_0$};
    \node[below] at (1.35,0) {$3f_0$};
    \node[below] at (2.25,0) {$5f_0$};
    \node[below] at (3.15,0) {$7f_0$};
  \end{scope}
}
\end{tikzpicture}

\caption{Ideal harmonic amplitudes relative to each waveform's fundamental.
Stems indicate discrete partials, not the broadened peaks of a finite FFT.
Real oscillators may intentionally depart from these shapes.}
\end{figure}

\paragraph{What to learn}
Peak spacing reveals harmonic structure; relative heights help explain
brightness. Equal peak voltage at the oscillator outputs does not guarantee
equal fundamental height or perceived loudness. To compare only spectral
shape, temporarily adjust analyzer input gains to align the fundamentals.
Return to equal gains before making a level comparison. A missing peak may
also lie below the selected Y Scale or be hidden beneath another trace.

\clearpage
\section{Comparing A Filter}\label{sec:filter}

Use two inputs to see what a processor changes while keeping the source
identical. A low-pass filter and a steady saw or noise source are enough.

\begin{figure}[!htp]
\centering
% Functional routing diagram, not a Rack patch screenshot.
\begin{tikzpicture}[x=1cm,y=1cm,font=\small,
  block/.style={draw=panelRule,rounded corners=2pt,fill=panelPaper,
    minimum width=2.3cm,minimum height=.8cm,align=center},
  wire/.style={-{Stealth[length=5pt]},line width=.9pt,draw=panelInk}]
  \node[block] (source) at (1.3,2) {Saw / noise};
  \node[block] (filter) at (6,2) {Filter};
  \node[block] (listen) at (11,2) {Listening path};
  \node[block,draw=channelOne] (red) at (3.6,0) {Fourier: red};
  \node[block,draw=channelTwo] (green) at (8.4,0) {Fourier: green};
  \draw[wire] (source)--(filter);
  \draw[wire] (filter)--(listen);
  \fill[panelInk] (3.6,2) circle[radius=.06];
  \fill[panelInk] (8.4,2) circle[radius=.06];
  \draw[wire,draw=channelOne] (3.6,2)--(red);
  \draw[wire,draw=channelTwo] (8.4,2)--(green);
  \node[above,text=panelMuted] at (3.6,2.1) {Before};
  \node[above,text=panelMuted] at (8.4,2.1) {After};
\end{tikzpicture}

\caption{Tap the signal before and after the filter. Analyzer cables are
branches of the listening path; Fourier is not an audio effect.}
\end{figure}

\paragraph{Starting point}
Load \textbf{Mastering}, then set Slope to zero and Smooth to None. Use
$0\,\mathrm{dB}$ gain on both connected inputs. Keep Length 4096 and
Hop $30\,\mathrm{ms}$. Choose Average $300\,\mathrm{ms}$ for noise, or
turn it off for a steady saw. Start with low filter resonance.

\begin{enumerate}
  \item Branch the source to the red input and the filter input. Patch
  the filter output to green and to your listening path.
  \item With the cutoff high, compare levels well below the cutoff. A
  constant difference here can be the filter's passband gain; it is not
  necessarily a change in spectral shape.
  \item Lower the cutoff slowly. The green trace should fall below red at
  higher frequencies. Raise resonance to look for a peak near the cutoff.
  Wait for Average to settle after each change.
  \item At a frequency with useful energy in both traces, place the cursor
  on each trace and subtract their dB levels. If red is $-12\,\mathrm{dB}$
  and green is $-24\,\mathrm{dB}$, the observed change is $-12\,\mathrm{dB}$:
  about one quarter of the amplitude at that frequency.
\end{enumerate}

\paragraph{Read the comparison carefully}
A saw samples the response only at its harmonics; noise fills more of the
spectrum but fluctuates. The ratio becomes unreliable where the source is
very weak. Fourier does not divide traces or measure phase automatically.
This is a visual magnitude comparison, not a complete transfer-function
measurement. Saturation, modulation, or other nonlinear behavior can create
new harmonics, so the result then depends on the source and its level.

\paragraph{Go further}
Use blue for a second filter or processing stage. Keep analyzer gains equal
and compare all taps simultaneously. Disable Filled Display if overlapping
areas obscure the curves. A processor's latency can make changing sounds
appear at different stages of their envelope; use a steady sound first.

\clearpage
\section{Separating Nearby Tones}\label{sec:resolution}

Two low notes can look like one broad peak. This experiment distinguishes
frequency resolution from simply enlarging the display.

\paragraph{Starting point}
Use two sine oscillators near $100$ and $110\,\mathrm{Hz}$, summed through
a mixer into red. At a Rack sample rate of $48\,\mathrm{kHz}$, choose Hann,
Length 2048, Hop $30\,\mathrm{ms}$, Linear Freq Scale, and
$0$--$300\,\mathrm{Hz}$. Set Average and Smooth off and Slope to zero.

\begin{enumerate}
  \item Observe the mixture, then mute each oscillator in turn to locate
  its contribution. At Length 2048 the two tones are closer than one bin.
  They will be difficult to distinguish in the mixture.
  \item Restore both tones and increase Length to 16384. Allow fresh audio
  to fill the new frame. The finer bin spacing should make them easier to
  separate with Hann; a broader window such as Flattop may still merge them.
  \item Return to Length 2048 and narrow LO/HI Freq around the peaks.
  The existing shape grows on screen, but the missing detail does not return.
  Turn off Bezier Curve in the context menu to see the bin connections.
  \item Return to Length 16384 and change one oscillator's pitch abruptly.
  Notice the transition taking time to pass through the long frame. A smaller
  Hop updates more frequently but does not shorten that frame.
\end{enumerate}

\begin{center}
\renewcommand{\arraystretch}{1.25}
\begin{tabular}{@{}rrr@{}}
\toprule
\textbf{Length at 48 kHz} & \textbf{Audio span} & \textbf{Bin spacing} \\
\midrule
2048 & $42.67\,\mathrm{ms}$ & $23.44\,\mathrm{Hz}$ \\
4096 & $85.33\,\mathrm{ms}$ & $11.72\,\mathrm{Hz}$ \\
8192 & $170.67\,\mathrm{ms}$ & $5.86\,\mathrm{Hz}$ \\
16384 & $341.33\,\mathrm{ms}$ & $2.93\,\mathrm{Hz}$ \\
\bottomrule
\end{tabular}
\end{center}

\paragraph{What sets the limit?}
Bin spacing is $f_s/N$; it is not a promise that every pair of tones that far
apart will resolve. Window width, relative levels, and leakage also matter.
Frequency smoothing trades more detail for a steadier broad shape. Increasing
Rack's sample rate at the same Length widens bin spacing; it does not improve
bass resolution. Choose Length for the detail needed, then Hop for how often
you need an update, while watching the patch's CPU load.

\clearpage
\section{Making Useful Measurements}\label{sec:measurements}

\paragraph{Establish a reference}
For a controlled check at $48\,\mathrm{kHz}$, use a $750\,\mathrm{Hz}$ sine:
with Length 2048 it falls exactly on bin 32. Use a known $5\,\mathrm{V}$
peak ($10\,\mathrm{V}$ peak-to-peak) source, unity input gain, Flattop,
zero Slope, no smoothing, and AC coupling off. After the frame fills,
the peak should be approximately $0\,\mathrm{dB}$ or $100\,\%$.
Use the source's frequency setting or a separate tuner to establish pitch;
the cursor does not snap to a peak. An off-bin tone can read differently.

\begin{center}
\renewcommand{\arraystretch}{1.25}
\begin{tabular}{@{}rrr@{}}
\toprule
\textbf{Sine peak voltage} & \textbf{Relative amplitude} & \textbf{Level} \\
\midrule
$5\,\mathrm{V}$ & $100\,\%$ & $0\,\mathrm{dB}$ \\
$2.5\,\mathrm{V}$ & $50\,\%$ & $-6.02\,\mathrm{dB}$ \\
$0.5\,\mathrm{V}$ & $10\,\%$ & $-20\,\mathrm{dB}$ \\
$0.005\,\mathrm{V}$ & $0.1\,\%$ & $-60\,\mathrm{dB}$ \\
\bottomrule
\end{tabular}
\end{center}

These ideal isolated-tone values follow $20\log_{10}(V_{\rm peak}/5\,\mathrm{V})$.
They are not a conversion from an arbitrary signal's largest spectral peak
to its total voltage. This reference is not dBFS, SPL, LUFS, or a calibrated
noise-power density measurement.

\paragraph{Know what is being summed}
Each jack sums all voices of a polyphonic cable before analysis. Two identical,
in-phase voices add to twice the amplitude (about $+6.02\,\mathrm{dB}$);
equal opposite-polarity voices cancel. To inspect voices independently, split
them into separate monophonic cables and use different inputs. Equal-looking
magnitude spectra alone cannot tell you the relative phase or predict the
result of mixing two signals.

\paragraph{Separate signal changes from viewing changes}
Input gain and AC coupling affect newly captured samples. Window, Average,
and Smooth affect the analysis. Frequency bounds and scales determine the
view; Slope deliberately weights levels by frequency. Turn Run off to compare
analysis choices on retained audio, but remember that Fourier keeps
reanalyzing and Average can still settle. Changing Length clears that audio.

\paragraph{Record the conditions}
When comparing a patch revision, note the Rack sample rate, input source and
gain, Length, Hop, window, smoothing, Slope, scales, and AC coupling. A noise
trace can change height when Length or window changes because each bin covers
a different effective bandwidth. Do not diagnose a quieter source from that
change alone. Save settings as a preset; save an external image if you need a
record of the trace, because patches do not retain captured spectra.
